\documentclass{article}

% Language setting
% Replace `english' with e.g. `spanish' to change the document language
\usepackage[english]{babel}

% Set page size and margins
% Replace `letterpaper' with`a4paper' for UK/EU standard size
\usepackage[letterpaper,top=2cm,bottom=2cm,left=3cm,right=3cm,marginparwidth=1.75cm]{geometry}

% Useful packages
\usepackage{amsmath}
\usepackage{amssymb}
\usepackage{amsthm}
\usepackage{amsfonts}
\usepackage{graphicx}
\usepackage{algorithm2e}
\usepackage{enumitem}
\usepackage{comment}
\usepackage{natbib}

\usepackage[colorlinks=true, allcolors=blue]{hyperref}

\newcommand{\sdag}[1]{{#1}^{\dag}}

\title{Fraction of Time MPI-SPPY spends in Python}
\author{ David L Woodruff\\
  Graduate School of Management\\
  \\
  University of California Davis\\
  Davis CA 95616 USA}
\date{\today}

\newtheorem{theorem}{Theorem}
\newtheorem{lemma}{Lemma}

\begin{document}
\maketitle

When considering applications in practice, or when comparing to other
packages, a question arises concerning the fraction of time that
mpi-sppy spends ``in Python'' as opposed to compiled code written in
other languages such as C and Fortran.  Since solvers, numpy, and MPI
are all in the latter category, {\em a priori} one expects that the
fraction spent in Python will be small.

We have used the tool called {\em scalene}, which is designed to
attribute time to Python vs native (it can estimate time spent in
compiled code called from Python, which it calls ``native''). It works
by sampling.

The short answer is: for hard problems, not much.
The longer answer is that there is no single number: the fraction
depends almost entirely on how much work the solver does per
subproblem. When the subproblems are even somewhat substantial, as in the
\texttt{sslp\_15\_45} instances, about 9\% of the time is spent in
Python, which matches the {\em a priori} expectation. For harder subproblems,
the fraction can be smaller. When the
subproblems are small enough that the solver returns almost
immediately, as in farmer, the majority of the time is spent in Python
--- 76\% in the most extreme case measured here. The crossover is not
subtle and it is not a sampling artifact.

The practical consequences are:

\begin{itemize}
  \item If your subproblems keep the solver busy, you cannot get much
    speed-up just by using another language. Maybe another language
    will make it easier for you to do your work.
  \item If your subproblems are small, the per-solve Python cost is
    what you are paying for, and the fix is to give the solver more
    work per call rather than to rewrite anything. Bundling is the
    obvious lever, and it is examined below.
\end{itemize}

\section{Method}

Every case runs the same three cylinders --- a PH hub, a lagrangian
spoke and an xhatshuffle spoke, one MPI rank each --- so that only the
model and the instance change between cases. The convergence tolerance
is set to zero so that runs end on the iteration limit rather than on a
gap, which keeps the amount of work per case predictable.

Because scalene samples, every case is run three times and the tables
report the mean over repetitions with the observed min--max range. The
spread turns out to be small for runs of a minute or so (typically a
few tenths of a percentage point) and much larger for the deliberately
short \texttt{farmer3} case, which is why that case is kept: it shows
what the numbers look like when there is not enough run time to sample
properly.

The numbers are read out of scalene's JSON profile rather than scraped
from the output of \texttt{scalene view --cli}. Summing the per-line
percentages in the JSON gives the same quantities at full precision,
and avoids three problems with scraping the terminal output: it rounds
each line to a whole percent, its \texttt{--reduced} form omits
low-usage lines so the sums undercount, and it is colourized, which
silently broke the row-matching in the earlier version of this code.
Scalene attributes 91--99\% of wall time to some source line; the
percentages below are shares of that attributed time, so the
unattributed remainder is divided out rather than being silently
credited to one of the buckets.

All runs are on a single machine with homogeneous cores, which avoids a
difficulty that affected an earlier version of these experiments: on a
machine with a mixture of fast and slow cores, ranks landing on slow
cores showed inflated Python fractions.

\noindent\textbf{System and software:}\\
\begin{itemize}
  \item \texttt{os}={Linux 7.0.0-28-generic (x86\_64)}
  \item \texttt{cpu\_model}={AMD Ryzen 7 7840HS}
  \item \texttt{cpu\_physical\_cores}={8}, \texttt{cpu\_logical}={16},
        \texttt{threads\_per\_core}={2}
  \item \texttt{cpu\_max\_mhz}={5137.9}
  \item \texttt{mem\_total}={14.31 GiB}
  \item \texttt{python}={3.11.13}, \texttt{scalene}={2.0.1},
        \texttt{pyomo}={6.9.5.dev0}, \texttt{mpi4py}={4.1.0},
        \texttt{gurobipy}={13.0.2}
  \item solver interface: \texttt{gurobi\_persistent},
        \texttt{--max-solver-threads 2}
\end{itemize}

\section{Results}

Table~\ref{tab:python-fraction-summary} gives the split by case and
Table~\ref{tab:python-fraction-by-rank} breaks the Python percentage
out by cylinder. The cases are:

\begin{itemize}
  \item \texttt{farmer3} --- 3 scenarios, 100 iterations. Deliberately
    short: about two seconds unprofiled, five under scalene.
  \item \texttt{farmer60}, \texttt{farmer240} --- 60 and 240
    scenarios, 400 and 200 iterations. Many trivially small LP
    subproblems.
  \item \texttt{farmer240\_bun10} --- \texttt{farmer240} with ten
    scenarios gathered into each proper bundle, so 24 subproblems
    instead of 240.
  \item \texttt{sslp\_15\_45\_10}, \texttt{sslp\_15\_45\_15} --- 10 and
    15 scenarios, 50 iterations. MIP subproblems that keep the solver
    busy.
  \item \texttt{sslp\_15\_45\_15\_bun3} ---
    \texttt{sslp\_15\_45\_15} with three scenarios per bundle, so 5
    subproblems instead of 15.
  \item \texttt{sslp\_5\_25\_50} --- 50 scenarios, 150 iterations. Many
    small MIP subproblems.
\end{itemize}

% Auto-generated by summarize_reps.py -- do not edit by hand

\begin{table}[ht]
\centering
\begin{tabular}{l r r r r r r}
\hline
Case & Reps & Wall (s) & Python (\%) & Native (\%) & System (\%) & Overhead \\
\hline
farmer3 & 3 & 5.3 & 16.2 (14.4--17.6) & 70.8 & 13.0 & 2.36$\times$ \\
farmer60 & 3 & 43.4 & 71.3 (70.3--72.4) & 16.2 & 12.4 & 1.55$\times$ \\
farmer240 & 3 & 82.9 & 76.1 (75.2--77.0) & 11.8 & 12.2 & 1.42$\times$ \\
farmer240\_bun10 & 3 & 20.9 & 62.6 (61.9--63.2) & 28.6 & 8.8 & 1.48$\times$ \\
sslp\_15\_45\_10 & 3 & 74.0 & 8.7 (8.5--9.0) & 89.3 & 2.0 & 1.09$\times$ \\
sslp\_15\_45\_15 & 3 & 104.8 & 9.6 (9.4--9.9) & 88.5 & 1.9 & 1.11$\times$ \\
sslp\_15\_45\_15\_bun3 & 3 & 73.7 & 12.5 (12.3--12.7) & 86.1 & 1.4 & 0.99$\times$ \\
sslp\_5\_25\_50 & 3 & 74.0 & 58.1 (58.0--58.3) & 37.1 & 4.7 & 1.22$\times$ \\
\hline
\end{tabular}
\caption{Time split by case, averaged over repetitions, with the min--max range over repetitions shown for the Python percentage. Percentages are of the time Scalene attributed to a source line (91.1--99.2\% of wall time here); wall time is the maximum over ranks. The overhead column is profiled wall time divided by unprofiled wall time for the same case; because scalene's instrumentation cost lands mostly on Python, it inflates the Python column.}
\label{tab:python-fraction-summary}
\end{table}

\begin{table}[ht]
\centering
\begin{tabular}{l r r r}
\hline
Case & PH hub & lagrangian & xhatshuffle \\
\hline
farmer3 & 14.1 (12.7--15.5) & 18.6 (17.5--20.5) & 15.9 (12.6--18.4) \\
farmer60 & 72.0 (71.3--72.9) & 72.6 (70.4--73.9) & 69.3 (68.6--70.4) \\
farmer240 & 75.9 (75.5--76.4) & 77.7 (76.6--78.6) & 74.6 (73.6--76.1) \\
farmer240\_bun10 & 62.3 (61.1--63.4) & 64.8 (64.1--65.7) & 60.8 (60.2--61.5) \\
sslp\_15\_45\_10 & 6.8 (6.5--7.3) & 6.3 (6.1--6.5) & 13.0 (12.8--13.2) \\
sslp\_15\_45\_15 & 7.2 (6.7--7.5) & 6.7 (6.3--7.0) & 14.9 (14.6--15.3) \\
sslp\_15\_45\_15\_bun3 & 5.1 (5.0--5.3) & 7.1 (6.9--7.2) & 25.3 (25.0--25.6) \\
sslp\_5\_25\_50 & 48.6 (48.4--48.8) & 48.8 (48.1--49.8) & 78.1 (77.2--79.2) \\
\hline
\end{tabular}
\caption{Python percentage by cylinder: mean over repetitions with the min--max range in parentheses.}
\label{tab:python-fraction-by-rank}
\end{table}


The ordering is by subproblem difficulty, not by instance size. The two
\texttt{sslp\_15\_45} cases sit at 8.7\% and 9.6\% Python; every case
whose subproblems are small sits far above that, from 58\% for
\texttt{sslp\_5\_25\_50} up to 76\% for \texttt{farmer240}. Note that
\texttt{sslp\_5\_25\_50} has more scenarios than
\texttt{sslp\_15\_45\_15} and yet spends six times the fraction of its
time in Python: what matters is the size of each subproblem, not how
many there are.

Looking at where the samples land makes the mechanism explicit. In both
regimes essentially all of the time is attributed to one line,
\texttt{spopt.py:337}, which is the Pyomo \texttt{solve} call. For
\texttt{sslp\_15\_45\_15} that line is 6.6\% Python and 88.2\% native
--- the solver is doing the work. For \texttt{farmer60} the same line
is 48.2\% Python and 5.8\% native, and a further 9.4\% of the run is
Python time in \texttt{spopt.py:267}, where the proximal objective is
handed to the solver on each iteration. That objective work scales with
the size of the model rather than with the difficulty of the solve,
which is exactly why it is invisible for sslp and dominant for farmer.

The per-cylinder table shows the PH hub and the lagrangian spoke
agreeing closely within every case, as they should, since they solve
the same subproblems with different objectives. The xhatshuffle spoke
is consistently the most Python-heavy of the three, and the gap widens
as the hub's own Python share falls: 6.8\% against 13.0\% for
\texttt{sslp\_15\_45\_10}, and 5.1\% against 25.3\% for
\texttt{sslp\_15\_45\_15\_bun3}. This is expected rather than
anomalous. The xhatshuffle spoke is looking for a feasible incumbent
rather than solving the subproblem to optimality, so much of its work
is Python bookkeeping --- shuffling the scenario order, fixing
candidate nonanticipative values, and checking the result --- with
comparatively little time left inside the solver. It is a small
fraction of total work, since it is one rank of three, but it is worth
knowing that the cylinders are not interchangeable for this
measurement.

\subsection{Bundles}

Bundling gives the solver more work per call, so it is the natural
response to a high Python fraction. It helps, but not reliably, and it
is worth being clear about why.

For farmer, bundling ten scenarios per subproblem takes the Python
share from 76.1\% down to 62.6\% and cuts wall time by a factor of
four, from 82.9 to 20.9 seconds. Fewer, larger solves means fewer
round trips through Pyomo, and the wall-clock gain is large. The Python
share nonetheless stays high, because a bundle of ten farmer scenarios
is still a trivial LP.

For \texttt{sslp\_15\_45\_15}, bundling three scenarios per subproblem
moves the Python share the other way, from 9.6\% up to 12.5\%, while
also cutting wall time from 104.8 to 73.7 seconds. Both numbers moved
in the direction that makes sense once the effect is separated into its
two parts: bundling reduced the total solver work more than it reduced
the total Python work, so Python's {\em share} rose even though the run
got faster.

So bundling should be judged on wall time, where it won in both cases,
and not on the Python fraction, which it can move either way. A high
Python fraction is a symptom worth investigating, but driving it down
is not itself the objective.

\subsection{Caveats}

Scalene's instrumentation is not free, and its cost falls mostly on
Python, so these Python percentages are upper bounds. The overhead
column in Table~\ref{tab:python-fraction-summary} is the ratio of
profiled to unprofiled wall time for the same case, and it lines up
with the Python column: the solver-bound sslp cases run at
0.99--1.11$\times$, essentially unaffected, while the Python-heavy
farmer cases run at 1.42--1.55$\times$. So the true Python fractions
for the farmer cases are somewhat lower than the table reports. The
exception is \texttt{farmer3} at 2.36$\times$, which is not a
Python-share effect at all: the profiler's fixed startup cost is simply
large next to a two-second run, which is another reason not to trust
that row for anything but its variability.

The gap between the two regimes is far too large to be an artifact of
this overhead --- 9\% against 76\% will not be closed by a factor of
1.5 --- but the individual percentages should be read as approximate.

Two further limits are worth stating. These are Pyomo models and the
Pyomo time is included in the Python total, so a good deal of what is
labelled Python here is Pyomo rather than mpi-sppy. And all of this is
one machine, one solver, and three repetitions per case.

\subsection{A note on non-persistent solvers}

All of the above uses \texttt{gurobi\_persistent}. If you use a
non-persistent interface instead --- \texttt{--solver-name gurobi} is
Pyomo's file-based interface, which writes an LP file and parses a
solution file on every solve --- then Pyomo does a great deal of extra
Python work per solve, which matters most when the subproblems are
small: \texttt{sslp\_5\_25\_50} goes from 58.1\% to 75.4\% Python and
from 74 to 124 seconds, and even the solver-bound
\texttt{sslp\_15\_45\_10}, whose wall time barely moves, goes from
8.7\% to 31.8\% Python. Use a persistent interface if you have one.

\end{document}
